\begin{frame}{왜 정밀도(Q1)가 첫 질문인가}
  \begin{block}{정밀도 = 필요충분 조건, 속도 = 비용 절감}
    정밀도(Q1)가 \textbf{필수 조건}이기 때문 --- 정밀도가 부족하면
    그 이상 속도가 빨라도 \textbf{잘못된 물리를 빠르게} 배우는
    것(13.1의 현실 격차가 그 출발점). 정밀도가 \textbf{충족된
    이후}에야 속도(Q2, Q3)가 ``어느 CPU/GPU''의 선택이 된다.
  \end{block}
  \vskip0.6em
  \begin{columns}[T]
    \begin{column}{0.48\textwidth}
      \kb{정밀도 (Q1)}
      \begin{itemize}
        \item \textbf{필수 조건} --- 부족하면 속도 무의미
        \item ``잘못된 물리를 빠르게 배우는'' 것을 막음
      \end{itemize}
    \end{column}
    \begin{column}{0.48\textwidth}
      \kb{속도 (Q2, Q3)}
      \begin{itemize}
        \item 정밀도 충족 \textbf{이후}의 선택
        \item ``어느 CPU/GPU''를 쓸지의 \textbf{비용 절감} 문제
      \end{itemize}
    \end{column}
  \end{columns}
\end{frame}
